\section{Geometric phase and gauge diagnostics}
\label{sec:geophase}

\subsection{Narrowband multi-component waves and polarization transport}

Assume the substrate supports a narrowband wave packet whose dominant oscillation can be expressed
locally as
\begin{equation}
  \vecu(x,t) \approx \mathrm{Re}\Big\{ a(x,t)\,\bm{\epsilon}(x,t)\,e^{i\theta(x,t)} \Big\},
  \qquad
  \theta(x,t)=\veck\cdot x-\omega t,
  \label{eq:quasi-monochromatic}
\end{equation}
where $\bm{\epsilon}(x,t)\in\mathbb{C}^N$ is a \emph{normalized} polarization state ($\bm{\epsilon}^\dagger\bm{\epsilon}=1$),
$a(x,t)\ge 0$ is a slowly varying amplitude, and $N$ counts the relevant internal components
(e.g., transverse directions in the embedding plus any effective internal basis used for coarse-graining).

The phase $\theta$ contains the fast oscillation; the \emph{geometric} information is encoded in how
$\bm{\epsilon}$ varies under transport in $(x,t)$.

\subsection{Berry / Pancharatnam--Berry connection}

Define the Berry connection one-form
\begin{equation}
  A_\mu := -i\,\bm{\epsilon}^\dagger \partial_\mu \bm{\epsilon},
  \qquad \mu\in\{t,1,2,3\},
  \label{eq:berry-connection}
\end{equation}
and its curvature
\begin{equation}
  F_{\mu\nu} := \partial_\mu A_\nu - \partial_\nu A_\mu.
  \label{eq:berry-curvature}
\end{equation}
Under a local phase change $\bm{\epsilon}\mapsto e^{i\chi(x,t)}\bm{\epsilon}$,
the connection transforms as $A_\mu\mapsto A_\mu+\partial_\mu \chi$ while $F_{\mu\nu}$ is invariant.
This is the standard $U(1)$ gauge freedom \citep{Berry1984,Pancharatnam1956,Cisowski2022}.

\paragraph{Gauge-invariant holonomy.}
For a closed loop $\gamma$ in configuration space,
\begin{equation}
  \Phi_\gamma := \oint_\gamma A_\mu\,\mathrm{d}x^\mu
\end{equation}
is gauge-invariant modulo $2\pi$ and equals the Berry phase accumulated along $\gamma$.
A key diagnostic is whether such loop phases behave \emph{field-like} (i.e., depend locally
on enclosed curvature and satisfy consistent evolution laws).

\subsection{Non-Abelian (Wilczek--Zee) connection for degenerate subspaces}

If the substrate supports a $k$-dimensional (near-)degenerate subspace of modes
$\{\bm{\epsilon}_a\}_{a=1}^k$ spanning a local vector bundle, define the matrix-valued connection
\begin{equation}
  (\mathcal{A}_\mu)_{ab} := -i\,\bm{\epsilon}_a^\dagger\,\partial_\mu \bm{\epsilon}_b,
  \label{eq:wz-connection}
\end{equation}
and curvature $\mathcal{F}_{\mu\nu}=\partial_\mu\mathcal{A}_\nu-\partial_\nu\mathcal{A}_\mu-i[\mathcal{A}_\mu,\mathcal{A}_\nu]$.
A change of local basis $\bm{\epsilon}_a\mapsto U_{ab}(x,t)\bm{\epsilon}_b$ with $U\in U(k)$ produces the
usual non-Abelian gauge transformation \citep{WilczekZee1984}.
This provides a principled way to encode multi-component ``internal'' structure without introducing
additional fundamental fields.

\subsection{Interpretation in the substrate program}

The model does \emph{not} assume that $A_\mu$ is electromagnetism. Instead:
\begin{itemize}
\item $A_\mu$ is a well-defined geometric object extracted from the substrate wave state.
\item $F_{\mu\nu}$ is a gauge-invariant diagnostic that can be tested for locality and propagation laws.
\item If, in regimes outside localized excitations, $F_{\mu\nu}$ obeys effective sourceless dynamics
and couples to localized modes through phase transport, then an emergent gauge interpretation becomes plausible.
\end{itemize}

The advantage of this viewpoint is methodological: geometric-phase quantities can be computed from
either continuum solutions or discrete simulations, providing a bridge from substrate to effective fields.
